Este Artigo insvestiga se a posição atual de um ponto de acesso(AP) Wi-Fi é adaquado para a cobertura de uma resisdencia. Para isso, foi definido uma planta baixa do ambiente e uma malha de pontos de medição, mantendo se o mesmo aparelho receptor, a mesma altura e condições equivalentes nas duas etapas. Inicialmente, foi registrado o nivel de potencia recebida em dBm com o ponto de acesso em uma posição original. Então, com base na distribuição observada e nas caracteristicas fisicas da residencia, foi formulada uma hipotese de reposicionamento. O AP foi movido do local original e a malha de medições foi repetida.
Os dados foram agregados pela mediana da leitura e realizadas em cada ponto e representados em dois mapas de calores comparaveis.
\medskip

\noindent \textbf{Palavras-chave:}  redes sem fio; Wi‐Fi; potência recebida; dBm; mapa de calor; posicionamento
de ponto de acesso \bigskip